\documentclass[dvisvgm,tikz,border=3pt]{standalone}
\usepackage{amsmath,amssymb}
\usepackage{pgfplots}
\usepackage{CJKutf8}
\pgfplotsset{compat=1.18}
\usetikzlibrary{arrows.meta,calc,positioning}

\definecolor{themebg}{HTML}{30362d}
\definecolor{themefg}{HTML}{edf4e8}
\definecolor{themegray}{HTML}{a4af9d}
\definecolor{themecurve}{HTML}{7eb6ff}
\definecolor{themealert}{HTML}{ff7b7b}
\definecolor{themeamber}{HTML}{f5b942}

\pgfdeclarelayer{background}
\pgfdeclarelayer{foreground}
\pgfsetlayers{background,main,foreground}

\begin{document}
\begin{CJK*}{UTF8}{gbsn}
\pagecolor{themebg}
\begin{tikzpicture}[color=themefg, text=themefg]

% ── 左图: 无穷区间反常积分 \int_1^\infty 1/x^p dx (要求 p > 1 衰减足够快) ──
\begin{axis}[
    name=axInf,
    width=7.8cm,
    height=7.2cm,
    axis lines=middle,
    axis line style={color=themefg, -{Stealth}},
    tick style={color=themefg},
    tick label style={color=themefg, font=\scriptsize},
    every axis label/.style={color=themefg},
    xmin=-0.2, xmax=5.2,
    ymin=-0.2, ymax=2.2,
    xtick={1},
    xticklabels={$1$},
    ytick={1},
    yticklabels={$1$},
    clip=false,
    xlabel={$x$},
    ylabel={$y$},
    title style={font=\footnotesize\bfseries, at={(0.5,1.06)}, anchor=south},
    title={无穷区间：$\int_1^{+\infty} \frac{1}{x^p} dx$},
]

% 积分尾部区域
\begin{pgfonlayer}{background}
    \addplot[fill=themecurve!30!themebg, domain=1:4.8, samples=50] {1/(x^1.5)} \closedcycle;
\end{pgfonlayer}

% 曲线对比
\addplot[themecurve, line width=2.0pt, domain=0.7:5.0, samples=60] {1/(x^1.5)};
\node[text=themecurve, font=\scriptsize\bfseries, anchor=south west] at (axis cs:2.2, 0.45) {$p>1$ (收敛, 衰减快)};

\addplot[themealert, dashed, line width=1.5pt, domain=0.7:5.0, samples=60] {1/x};
\node[text=themealert, font=\scriptsize\bfseries, anchor=south west] at (axis cs:2.8, 0.8) {$p=1$ (发散)};

\begin{pgfonlayer}{foreground}
    \node[text=themefg, font=\scriptsize, align=left, anchor=north west] at (axis cs:0.2, 2.05) {
        \textbf{无穷区间判敛准则}：\\
        $\boxed{\int_1^{+\infty}\frac{1}{x^p}dx\text{ 收敛} \iff p > 1}$\\
        本质：尾部面积无限延伸，\\
        要求被积函数衰减\textbf{足够快}
    };
\end{pgfonlayer}
\end{axis}

% ── 右图: 无界瑕点反常积分 \int_0^1 1/x^p dx (要求 p < 1 奇性足够弱) ──
\begin{axis}[
    name=axSing,
    at={($(axInf.east)+(1.6cm,0)$)},
    anchor=west,
    width=7.8cm,
    height=7.2cm,
    axis lines=middle,
    axis line style={color=themefg, -{Stealth}},
    tick style={color=themefg},
    tick label style={color=themefg, font=\scriptsize},
    every axis label/.style={color=themefg},
    xmin=-0.2, xmax=2.2,
    ymin=-0.2, ymax=4.2,
    xtick={1},
    xticklabels={$1$},
    ytick={1},
    yticklabels={$1$},
    clip=false,
    xlabel={$x$},
    ylabel={$y$},
    title style={font=\footnotesize\bfseries, at={(0.5,1.06)}, anchor=south},
    title={无界瑕点：$\int_0^1 \frac{1}{x^p} dx$},
]

% 瑕点积分区域
\begin{pgfonlayer}{background}
    \addplot[fill=themeamber!30!themebg, domain=0.06:1, samples=60] {1/sqrt(x)} \closedcycle;
\end{pgfonlayer}

% 曲线对比
\addplot[themeamber, line width=2.0pt, domain=0.06:1.8, samples=70] {1/sqrt(x)};
\node[text=themeamber, font=\scriptsize\bfseries, anchor=west] at (axis cs:0.4, 2.4) {$p<1$ (收敛, 奇性弱)};

\addplot[themealert, dashed, line width=1.5pt, domain=0.25:1.8, samples=50] {1/x};
\node[text=themealert, font=\scriptsize\bfseries, anchor=west] at (axis cs:0.85, 3.4) {$p=1$ (发散)};

\begin{pgfonlayer}{foreground}
    \node[text=themefg, font=\scriptsize, align=left, anchor=north east] at (axis cs:2.1, 4.1) {
        \textbf{无界瑕点判敛准则}：\\
        $\boxed{\int_0^1\frac{1}{x^p}dx\text{ 收敛} \iff p < 1}$\\
        本质：原点垂向无限冲高，\\
        要求被积函数奇性\textbf{足够弱}
    };
\end{pgfonlayer}
\end{axis}

\end{tikzpicture}
\end{CJK*}
\end{document}
